% GDZ PPN236018647,native594lower–595through133;HKUST595–596.
\NumberedItem{130}
\noindent\hspace*{\parindent}Theoria Residuorum cubicorum et biquadraticorum incepta
\par\noindent\hfill 1807 Febr. 15,

\NumberedItem{131}
\noindent ulterius exculta et completa reddita Febr. 17. Demonstratione adhuc eget:
\par\noindent\hfill [1807 Febr. 17.]
\clearpage\SourceLeaf{595}\ArithmeticHead{566}{L}{TAGEBUCH MIT ERLÄUTERUNGEN.}
\NumberedItem{132}
\noindent\hspace*{\parindent}Demonstratio huius theoriae per methodum elegantissimam inventa ita ut penitus perfecta sit nihilque amplius desideretur.
\par\noindent\hspace*{\parindent}Hinc simul residua et non residua quadratica egregie illustrantur.
\par\noindent\hfill 1807 Febr. 22.

\NumberedItem{133}
\noindent\hspace*{\parindent}Theoremata, quae theoriae praecedenti incrementa maximi pretii adiungunt, demonstratione eleganti munita (scilicet pro quibusnam radicibus primitivis statuere oporteat ipsum \(b\) positivum pro quibusque negativum,
\[
aa+27bb=4p;\quad aa+4bb=p).
\]
\par\noindent\hfill [1807] Febr. 24.
\par\medskip\begingroup\small
\noindent\hspace*{\parindent}In den Nummern 130—133 sind die Ergebnisse der \textit{Theoria Residuorum biquadraticorum, Commentatio prima} (vorgelegt am 5. April 1825), Werke II, S. 65 gemeint. In dieser Abhandlung (a. a. O. S. 67) und in ihrer Anzeige (ebenda S. 165), ebenso in dem Briefe an \textsc{Dirichlet} vom 30. Mai 1828 (ebenda, zweiter Abdruck, S. 516) setzt \textsc{Gauss} den Anfang seiner Beschäftigung mit den biquadratischen Resten auf 1805; in dem Briefe an \textsc{Sophie Germain} vom 30. April 1807 (siehe oben S. 70, insbesondere S. 72) wird »der letzte Winter« — also die Zeit der hier vorliegenden Tagebuchaufzeichnungen erwähnt. Vergl. auch die Bemerkung zu der Nr. 128, sowie die Werke VIII, S. 3—11 und 15—19 abgedruckten Nachlaßstücke und die zugehörigen Bemerkungen, ferner die oben S. 37 abgedruckte Aufzeichnung und die Briefstellen an \textsc{Sophie Germain} (oben S. 72, I, II) und an \textsc{Olbers} (oben S. 75).
\par\noindent\hfill\textsc{Klein.\quad Schlesinger.}
\par\endgroup
